\documentclass[11pt,a4paper]{article}

% ── Packages ──
\usepackage[utf8]{inputenc}
\usepackage[T1]{fontenc}
\usepackage{amsmath,amssymb}
\usepackage{physics}
\usepackage{siunitx}
\usepackage{booktabs}
\usepackage{geometry}
\usepackage{hyperref}
\usepackage{enumitem}
\usepackage{float}

\geometry{margin=2.5cm}
\hypersetup{colorlinks=true,linkcolor=blue,citecolor=blue,urlcolor=blue}

% ── Load computed values ──
% MARS parameters - auto-generated, do not edit
% Generated by mars_latex.py on 2026-03-22 21:29:49

% Fixed parameters
\newcommand{\Pbar}{15}  % bar
\newcommand{\TK}{300}  % K
\newcommand{\gammaXe}{1.667}  % 
\newcommand{\Rcyl}{1.6}  % m
\newcommand{\Hcyl}{1.5}  % m
\newcommand{\Ldrift}{1.5}  % m
\newcommand{\vdrift}{100}  % mm/s
\newcommand{\sigmadiff}{1}  % mm
\newcommand{\uc}{0.1}  % mm/s

% Design parameters
\newcommand{\narms}{2}  % 
\newcommand{\Rblade}{1.6}  % m
\newcommand{\chord}{0.16}  % m
\newcommand{\chordmm}{160}  % mm
\newcommand{\mblade}{0.25}  % kg
\newcommand{\mplate}{0.25}  % kg
\newcommand{\taumotor}{60}  % N.m
\newcommand{\deltar}{0.8}  % m
\newcommand{\rzerocarrier}{0.8}  % m
\newcommand{\zlift}{4}  % mm
\newcommand{\tlift}{0.4}  % s
\newcommand{\Factuator}{20}  % N

% Gas properties
\newcommand{\rhoxe}{79}  % kg/m3
\newcommand{\nuxe}{2.94e-07}  % m2/s
\newcommand{\csound}{178}  % m/s

% Phase 1: Blade rotation
\newcommand{\sweepangle}{180}  % deg
\newcommand{\trot}{395}  % ms
\newcommand{\trots}{0.395}  % s
\newcommand{\omegamax}{15.9}  % rad/s
\newcommand{\Vtip}{25.4}  % m/s
\newcommand{\Matip}{0.14}  % 
\newcommand{\Retip}{1.38e+07}  % 
\newcommand{\deltablade}{2.21}  % mm
\newcommand{\deltastar}{0.276}  % mm
\newcommand{\uprimeblade}{1.27}  % m/s
\newcommand{\maxreachblade}{2.9}  % mm
\newcommand{\tdecayblade}{3.5}  % ms

% Phase 2: Carrier slide
\newcommand{\tslide}{200}  % ms
\newcommand{\tslides}{0.2}  % s
\newcommand{\Vslide}{8}  % m/s
\newcommand{\deltacarrier}{2.78}  % mm
\newcommand{\maxreachcarrier}{3.68}  % mm
\newcommand{\tdecaycarrier}{13.9}  % ms

% Phase 3: Vertical lift
\newcommand{\Vlift}{20}  % mm/s
\newcommand{\Relift}{272}  % 
\newcommand{\deltaStokes}{0.34}  % mm

% Summary
\newcommand{\ttotal}{1}  % s
\newcommand{\ttotalms}{995}  % ms
\newcommand{\maxreach}{3.68}  % mm
\newcommand{\margin}{0.32}  % mm
\newcommand{\deadzone}{100}  % mm
\newcommand{\deadzonepct}{6.6}  % \%

% 6 Mechanisms
\newcommand{\tauDrag}{63}  % N.m
\newcommand{\deltaEkman}{0.19}  % mm
\newcommand{\tacoustic}{9}  % ms
\newcommand{\dpPulse}{51}  % kPa
\newcommand{\acent}{101}  % m/s2
\newcommand{\fshed}{278}  % Hz


\title{MARS Flow Analysis for the ITACA Detector}
\author{}
\date{}

\begin{document}
\maketitle

\begin{abstract}
The MARS (Mechanically Actuated Rotor System) positions ion collection plates
in the ITACA detector, a high-pressure xenon chamber for ionization track
measurement. We analyze all gas flow mechanisms induced by the three-phase
MARS motion sequence: blade rotation, carrier slide, and vertical lift.
The velocity field is characterized by three components: $u_\theta(z)$ (tangential),
$u_r(z)$ (radial), and $u_z(z)$ (vertical), where $z$ is the height above the
blade/carrier surface. All six identified mechanisms have \textbf{zero effect
on ion collection}: horizontal velocities ($u_\theta$, $u_r$) are geometrically
separated by the lift, and transient disturbances decay in milliseconds while
ions arrive seconds later.
\end{abstract}


% ══════════════════════════════════════════════════════════════════════════════
\section{The ITACA Detector}
\label{sec:itaca}
% ══════════════════════════════════════════════════════════════════════════════

ITACA (Ion Tracking And Collection Apparatus) measures ionization tracks in
dense xenon gas. An ionization event produces electron-ion pairs along a track.
Electrons drift rapidly to the anode ($\sim$1~ms), providing the track position
$(r, \theta)$. Ions drift slowly downward and must be collected on a plate
positioned at the correct location.

\subsection{Chamber Geometry}

\begin{table}[H]
\centering
\caption{ITACA chamber parameters.}
\begin{tabular}{lll}
\toprule
Parameter & Value & Notes \\
\midrule
Shape & Vertical cylinder & \\
Diameter & \SI{3.2}{m} & $R_\mathrm{cyl} = \Rcyl$~m \\
Height & $H = \Hcyl$~m & Drift length \\
Gas & Xenon & \\
Pressure & $P = \Pbar$~bar & \\
Temperature & $T = \TK$~K & \\
\bottomrule
\end{tabular}
\end{table}

\subsection{Coordinate System}

We use cylindrical coordinates $(r, \theta, z)$:
\begin{itemize}[nosep]
  \item $z$-axis: Vertical (along cylinder axis), positive upward
  \item $z = H$ (\Hcyl~m): Anode (top)
  \item $z = 0$: Blade/carrier surface (before lift)
  \item $z = z_\mathrm{lift}$: Ion plate surface (after lift)
  \item $r$: Radial distance from central axis
  \item $\theta$: Azimuthal angle
\end{itemize}

The velocity field has three components:
\begin{itemize}[nosep]
  \item $u_\theta(z)$: Tangential velocity (perpendicular to radius, horizontal)
  \item $u_r(z)$: Radial velocity (along radius, horizontal)
  \item $u_z(z)$: Vertical velocity (along $z$-axis)
\end{itemize}

\subsection{Electrode Structure}

From top to bottom:
\begin{enumerate}[nosep]
  \item \textbf{Anode} at $z = H$: Collects electrons, provides $(r,\theta)$ readout
  \item \textbf{Drift region}: Ions produced here, drift downward at $v_d = \vdrift$~mm/s
  \item \textbf{Cathode grid} at $z \approx 10$~mm: Wire mesh, 92\% transparent
  \item \textbf{Ion plate} at $z = \zlift$~mm (after lift): Collects ions
  \item \textbf{MARS surface} at $z = 0$: Blade and carrier before lift
\end{enumerate}

\subsection{Ion Physics}

Ions are created throughout the drift region and travel downward at velocity $v_d$.
The minimum drift time corresponds to ions created just above the cathode grid
($z_\mathrm{cathode} \approx 10$~mm):
\begin{equation}
  t_\mathrm{drift,min} = \frac{z_\mathrm{cathode}}{v_d} = \frac{10~\text{mm}}{100~\text{mm/s}} = 0.1~\text{s}
\end{equation}

However, this assumes ions exactly at the cathode grid. In practice, a small
buffer ($z_\mathrm{buffer} \approx 100$~mm) is needed, giving:
\begin{equation}
  t_\mathrm{available} \approx \frac{z_\mathrm{buffer}}{v_d} = \frac{100~\text{mm}}{100~\text{mm/s}} = 1~\text{s}
\end{equation}

\begin{table}[H]
\centering
\caption{Ion drift parameters.}
\begin{tabular}{lll}
\toprule
Parameter & Value & Notes \\
\midrule
Drift velocity & $v_d = \vdrift$~mm/s & Downward ($-z$) \\
Cathode grid height & $z_\mathrm{cathode} \approx 10$~mm & Above MARS surface \\
Dead zone & $z_\mathrm{dead} = v_d \times t_\mathrm{total}$ & Ions lost during motion \\
Available time & $t_\mathrm{available} \sim 1$~s & Before first collectible ions arrive \\
Maximum drift time & $H/v_d \sim 15$~s & Ions near anode \\
Diffusion (RMS) & $\sigma_\mathrm{diff} \approx \sigmadiff$~mm & At cathode \\
\bottomrule
\end{tabular}
\end{table}

The electron signal at $t = 0$ triggers MARS. MARS must complete its motion
within $t_\mathrm{available}$ to collect ions from the useful drift region.


% ══════════════════════════════════════════════════════════════════════════════
\section{The MARS System}
\label{sec:mars}
% ══════════════════════════════════════════════════════════════════════════════

MARS positions the ion collection plate directly under the ionization track
location $(r, \theta)$ provided by the anode readout.

\subsection{Components}

\begin{enumerate}[nosep]
  \item \textbf{Propeller blades}: $n = \narms$ NACA\,0012 airfoils rotating about central axis
  \item \textbf{Carrier plates}: Ion collection plates mounted on blades
  \item \textbf{Linear actuators}: Slide carriers radially along blades
  \item \textbf{Lift mechanism}: Raise carrier plate vertically toward cathode
\end{enumerate}

\subsection{Blade Specification}

\begin{table}[H]
\centering
\caption{Blade parameters.}
\begin{tabular}{lll}
\toprule
Parameter & Value & Notes \\
\midrule
Number of arms & $n = \narms$ & \\
Length (radius) & $R = \Rblade$~m & Hub to tip \\
Chord & $c = \chordmm$~mm & \\
Profile & NACA\,0012 & Symmetric, 12\% thickness \\
Angle of attack & $\alpha = 0°$ & No lift generated \\
Mass per blade & $m_\mathrm{blade} = \mblade$~kg & \\
Motor torque & $\tau = \taumotor$~N$\cdot$m & \\
\bottomrule
\end{tabular}
\end{table}

\subsection{Carrier Plate Specification}

\begin{table}[H]
\centering
\caption{Carrier plate parameters.}
\begin{tabular}{lll}
\toprule
Parameter & Value & Notes \\
\midrule
Dimensions & $160 \times 160 \times 5$~mm$^3$ & \\
Mass & $m_\mathrm{plate} = \mplate$~kg & \\
Initial position & $r_0 = \rzerocarrier$~m & $= R/2$ \\
Slide distance & $\Delta r = \deltar$~m & \\
Actuator force & $F = \Factuator$~N & \\
\bottomrule
\end{tabular}
\end{table}

\subsection{Three-Phase Motion Sequence}

All motions are \textbf{sequential} (not simultaneous), triggered at $t = 0$
by the anode readout:

\begin{enumerate}[nosep]
  \item \textbf{Phase 1: Blade rotation} --- Tangential motion ($\theta$-direction)
  \item \textbf{Phase 2: Carrier slide} --- Radial motion ($r$-direction)
  \item \textbf{Phase 3: Vertical lift} --- Axial motion ($z$-direction)
\end{enumerate}

Each phase generates a distinct velocity component in the gas.


% ══════════════════════════════════════════════════════════════════════════════
\section{Xenon Gas Properties}
\label{sec:gas}
% ══════════════════════════════════════════════════════════════════════════════

Dense xenon at \Pbar~bar, \TK~K has unusual properties:

\begin{table}[H]
\centering
\caption{Xenon gas properties.}
\begin{tabular}{llll}
\toprule
Property & Symbol & Value & vs.\ air \\
\midrule
Density & $\rho$ & \rhoxe~kg/m$^3$ & 66$\times$ denser \\
Dynamic viscosity & $\mu$ & $2.32 \times 10^{-5}$~Pa$\cdot$s & Similar \\
Kinematic viscosity & $\nu$ & $\nuxe$~m$^2$/s & 51$\times$ smaller \\
Speed of sound & $c_s$ & \csound~m/s & --- \\
Heat capacity ratio & $\gamma$ & 5/3 & Monatomic \\
\bottomrule
\end{tabular}
\end{table}

\textbf{Key consequences:}
\begin{itemize}[nosep]
  \item High Reynolds numbers at modest velocities (turbulent flow)
  \item Very slow molecular diffusion ($t_\mathrm{diff} \sim \delta^2/\nu \sim$ hours)
  \item Long viscous timescales
\end{itemize}


% ══════════════════════════════════════════════════════════════════════════════
\section{Boundary Layer Theory}
\label{sec:bl-theory}
% ══════════════════════════════════════════════════════════════════════════════

\subsection{Flat-Plate Approximation}

We model the boundary layers on the blade and carrier using the
\textbf{flat-plate approximation}. This is justified because:
\begin{itemize}[nosep]
  \item The NACA\,0012 profile is thin (12\% thickness) and symmetric
  \item At $\alpha = 0°$, there is no pressure gradient along the chord
  \item The carrier plate is geometrically flat
\end{itemize}

For a turbulent boundary layer over a flat plate of length $L$ at Reynolds
number $\mathrm{Re}_L = V L / \nu$, the Prandtl--Schlichting correlation gives
the boundary layer thickness:
\begin{equation}
  \boxed{\delta = 0.37 \, L \, \mathrm{Re}_L^{-1/5}}
  \label{eq:delta-flat-plate}
\end{equation}

This formula is valid for $5 \times 10^5 < \mathrm{Re}_L < 10^7$ (fully turbulent).

\subsection{The 1/7 Power-Law Velocity Profile}

Within the turbulent boundary layer, the mean velocity profile follows the
empirical 1/7 power law:
\begin{equation}
  \boxed{u(z) = V \left(\frac{z}{\delta}\right)^{1/7} \quad \text{for } z < \delta}
  \label{eq:one-seventh}
\end{equation}
where $V$ is the freestream velocity and $z$ is the distance from the surface.

This profile:
\begin{itemize}[nosep]
  \item Matches experimental data for turbulent flat-plate BLs
  \item Has $u(0) = 0$ (no-slip at wall)
  \item Has $u(\delta) = V$ (matches freestream at BL edge)
  \item Has a steep gradient near the wall: $u(0.1\delta) \approx 0.72\,V$
\end{itemize}

\subsection{Displacement and Momentum Thickness}

For the 1/7 power-law profile, the integral thicknesses are:
\begin{align}
  \delta^* &= \int_0^\delta \left(1 - \frac{u}{V}\right) dz = \frac{\delta}{8}
    \label{eq:delta-star} \\
  \theta &= \int_0^\delta \frac{u}{V}\left(1 - \frac{u}{V}\right) dz = \frac{7\delta}{72}
    \label{eq:theta}
\end{align}

\subsection{Turbulence Intensity and Eddy Timescale}

The turbulence intensity in a flat-plate boundary layer is approximately:
\begin{equation}
  u' \approx 0.05 \, V
  \label{eq:u-prime}
\end{equation}
where $u'$ is the RMS velocity fluctuation.

The characteristic eddy turnover time is:
\begin{equation}
  t_\mathrm{eddy} = \frac{\delta}{u'}
  \label{eq:t-eddy}
\end{equation}

After the surface stops moving, turbulent eddies decay. The decay time is
approximately:
\begin{equation}
  \boxed{t_\mathrm{decay} = 2 \, t_\mathrm{eddy} = \frac{2\delta}{u'} = \frac{2\delta}{0.05\,V} = \frac{40\,\delta}{V}}
  \label{eq:t-decay}
\end{equation}

\subsection{Aerodynamic Drag}

The equation of motion for a rotating blade including drag is:
\begin{equation}
  I \frac{d\omega}{dt} = \tau_\mathrm{motor} - \tau_\mathrm{drag}
  \label{eq:eom-drag}
\end{equation}
where the drag torque is:
\begin{equation}
  \tau_\mathrm{drag} = \frac{1}{2} \rho \, C_d \, c \int_0^R (\omega r)^2 \, r \, dr
    = \frac{1}{8} \rho \, C_d \, c \, R^4 \, \omega^2
  \label{eq:tau-drag}
\end{equation}

Since $\tau_\mathrm{drag} \propto \omega^2$, the equation becomes nonlinear and
requires numerical integration. The analytical bang-bang solution assumes
$\tau_\mathrm{drag} = 0$.

\textbf{Drag coefficients:}
\begin{itemize}[nosep]
  \item NACA\,0012 at $\alpha = 0°$: $C_d \approx 0.006$--$0.01$ (highly streamlined)
  \item Carrier plate (sharpened edges): $C_d \approx 0.1$ (bluff body, reduced by edge treatment)
\end{itemize}

Numerical integration shows that including drag increases motion times by $<5\%$.
The analytical approximation (neglecting drag) is therefore acceptable for design purposes.

\subsection{Boundary Layer Scaling}

The Prandtl--Schlichting formula can be rewritten as:
\begin{equation}
  \delta = 0.37 \, L^{4/5} \left(\frac{V}{\nu}\right)^{-1/5}
  \quad \Rightarrow \quad
  \delta \propto L^{4/5} \, V^{-1/5}
  \label{eq:delta-scaling}
\end{equation}

For the blade (NACA) and carrier, the characteristic lengths are both $L = 160$~mm
(chord and plate width). The ratio of boundary layer thicknesses is:
\begin{equation}
  \frac{\delta_\mathrm{carrier}}{\delta_\mathrm{blade}}
    = \left(\frac{V_\mathrm{tip}}{V_\mathrm{slide}}\right)^{1/5}
    = \left(\frac{25.4}{8}\right)^{0.2}
    = 1.26
\end{equation}

This explains why $\delta_\mathrm{blade} = \deltablade$~mm and
$\delta_\mathrm{carrier} = \deltacarrier$~mm are similar despite the 3$\times$
velocity difference: the weak $V^{-1/5}$ scaling means a factor of 3 in velocity
produces only a 26\% change in $\delta$.


% ══════════════════════════════════════════════════════════════════════════════
\section{Three-Phase Motion Analysis}
\label{sec:phases}
% ══════════════════════════════════════════════════════════════════════════════

\subsection{Phase 1: Blade Rotation --- $u_\theta(z)$}

The propeller executes a bang-bang motion (maximum acceleration then deceleration)
through sweep angle $\theta_\mathrm{sweep} = 360°/n = \sweepangle°$.

\subsubsection{Kinematics}

\begin{align}
  I_\mathrm{total} &= n \left(\frac{1}{3}m_\mathrm{blade}R^2 + m_\mathrm{plate}r_0^2\right) \\
  t_\mathrm{rot} &= 2\sqrt{\frac{\theta_\mathrm{sweep} \cdot I_\mathrm{total}}{\tau}} = \trot~\text{ms} \\
  \omega_\mathrm{max} &= \sqrt{\frac{\theta_\mathrm{sweep} \cdot \tau}{I_\mathrm{total}}} = \omegamax~\text{rad/s} \\
  V_\mathrm{tip} &= \omega_\mathrm{max} R = \Vtip~\text{m/s}
\end{align}

\subsubsection{Flow Regime}

\begin{align}
  \mathrm{Ma} &= \frac{V_\mathrm{tip}}{c_s} = \Matip \quad \text{(incompressible, Ma} < 0.3\text{)} \\
  \mathrm{Re}_c &= \frac{V_\mathrm{tip} \cdot c}{\nu} = \Retip \quad \text{(fully turbulent)}
\end{align}

\subsubsection{Blade Boundary Layer}

Using the flat-plate correlation \eqref{eq:delta-flat-plate} with $L = c$:
\begin{equation}
  \delta_\mathrm{blade} = 0.37 \, c \, \mathrm{Re}_c^{-1/5} = \deltablade~\text{mm}
\end{equation}

The velocity profile within this boundary layer is \textbf{tangential}:
\begin{equation}
  \boxed{u_\theta(z) = V_\mathrm{tip} \left(\frac{z}{\delta_\mathrm{blade}}\right)^{1/7}
    \quad \text{for } z < \delta_\mathrm{blade}}
  \label{eq:u-theta}
\end{equation}

This is a \textbf{horizontal} velocity (perpendicular to radius).

\subsubsection{Turbulence Parameters}

\begin{align}
  u'_\mathrm{blade} &= 0.05 \times V_\mathrm{tip} = \uprimeblade~\text{m/s} \\
  t_\mathrm{decay,blade} &= \frac{2\,\delta_\mathrm{blade}}{u'_\mathrm{blade}} = \tdecayblade~\text{ms}
\end{align}

\subsubsection{Max Reach}

After the blade stops, turbulent eddies transport horizontal momentum vertically
during their lifetime $t_\mathrm{decay}$. The maximum height reached is:
\begin{equation}
  \mathrm{max\_reach}_\mathrm{blade} = \delta_\mathrm{blade} + \ell_\mathrm{turb} + \ell_\mathrm{mol}
    = \maxreachblade~\text{mm}
\end{equation}
where $\ell_\mathrm{turb} = \sqrt{\nu_t \, t_\mathrm{decay}}$ is the turbulent
diffusion length and $\ell_\mathrm{mol} = \sqrt{\nu \, t_\mathrm{decay}}$ is
the molecular diffusion length.

\subsection{Phase 2: Carrier Slide --- $u_r(z)$}

The carrier plate slides radially inward by $\Delta r = \deltar$~m using a
bang-bang profile.

\subsubsection{Kinematics}

\begin{align}
  t_\mathrm{slide} &= 2\sqrt{\frac{\Delta r \cdot m_\mathrm{plate}}{F}} = \tslide~\text{ms} \\
  V_\mathrm{slide} &= \sqrt{\frac{\Delta r \cdot F}{m_\mathrm{plate}}} = \Vslide~\text{m/s}
\end{align}

\subsubsection{Carrier Boundary Layer}

Using the flat-plate correlation with $L = $ carrier width (160~mm):
\begin{equation}
  \delta_\mathrm{carrier} = 0.37 \times 0.16 \times \mathrm{Re}_\mathrm{carrier}^{-1/5}
    = \deltacarrier~\text{mm}
\end{equation}

The velocity profile within this boundary layer is \textbf{radial}:
\begin{equation}
  \boxed{u_r(z) = V_\mathrm{slide} \left(\frac{z}{\delta_\mathrm{carrier}}\right)^{1/7}
    \quad \text{for } z < \delta_\mathrm{carrier}}
  \label{eq:u-r}
\end{equation}

This is a \textbf{horizontal} velocity (along the radius).

\subsubsection{Max Reach}

\begin{equation}
  \mathrm{max\_reach}_\mathrm{carrier} = \maxreachcarrier~\text{mm}
\end{equation}

\subsection{Phase 3: Vertical Lift --- $u_z(z)$}

The ion plate lifts vertically by $z_\mathrm{lift} = \zlift$~mm in
$t_\mathrm{lift} = \tlift$~s.

\subsubsection{Kinematics}

\begin{equation}
  V_\mathrm{lift} = \frac{2 z_\mathrm{lift}}{t_\mathrm{lift}} = \Vlift~\text{mm/s}
    \quad \text{(bang-bang peak)}
\end{equation}

\subsubsection{Flow Regime}

\begin{equation}
  \mathrm{Re}_\mathrm{lift} = \frac{V_\mathrm{lift} \cdot z_\mathrm{lift}}{\nu}
    = \Relift \quad \text{(laminar, Re} < 2000\text{)}
\end{equation}

\subsubsection{Stokes Layer}

The lift creates a laminar Stokes diffusion layer:
\begin{equation}
  \delta_\mathrm{Stokes} = \sqrt{\nu \cdot t_\mathrm{lift}} = \deltaStokes~\text{mm}
\end{equation}

The velocity profile decays exponentially with height:
\begin{equation}
  \boxed{u_z(z) = V_\mathrm{lift} \exp\left(-\frac{z}{\delta_\mathrm{Stokes}}\right)}
  \label{eq:u-z}
\end{equation}

This is a \textbf{vertical} velocity. It does not displace ions laterally---only
a small timing shift of $\sim$1~ms (negligible on 1--15~s drift).

\subsection{Summary: Velocity Components}

\begin{table}[H]
\centering
\caption{Velocity components from each phase.}
\begin{tabular}{llllll}
\toprule
Phase & Motion & Component & Direction & Formula & Effect on ions \\
\midrule
1 & Blade rotation & $u_\theta(z)$ & Tangential & Eq.~\eqref{eq:u-theta} & Lateral displacement \\
2 & Carrier slide & $u_r(z)$ & Radial & Eq.~\eqref{eq:u-r} & Lateral displacement \\
3 & Vertical lift & $u_z(z)$ & Vertical & Eq.~\eqref{eq:u-z} & Timing shift only \\
\bottomrule
\end{tabular}
\end{table}

\textbf{Key insight:} Only horizontal velocities ($u_\theta$, $u_r$) can displace
ions laterally. The vertical velocity $u_z$ only affects ion arrival timing.

\subsection{Velocity Profile Table}

Using the 1/7 power law for $u_\theta$ and $u_r$, and exponential decay for $u_z$:

\begin{table}[H]
\centering
\caption{Velocity profiles vs.\ height $z$ for $n = 2$ configuration.}
\begin{tabular}{rrrrll}
\toprule
$z$ (mm) & $u_\theta$ (m/s) & $u_r$ (m/s) & $u_z$ (mm/s) & Direction & Status \\
\midrule
0.1 & 16.3 & 5.0 & 15.4 & tan + rad + vert & Inside both BLs \\
0.5 & 20.6 & 6.3 & 5.4 & tan + rad + vert & Inside both BLs \\
1.0 & 22.7 & 6.9 & 1.5 & tan + rad + vert & Inside both BLs \\
2.0 & 25.1 & 7.6 & 0.1 & tan + rad & Inside both BLs \\
2.5 & 0 & 7.9 & 0 & rad only & Above blade BL \\
3.0 & 0 & 0 & 0 & --- & Quiescent \\
4.0 & 0 & 0 & 0 & --- & Quiescent \\
5.0 & 0 & 0 & 0 & --- & Ion plate surface \\
\bottomrule
\end{tabular}
\end{table}


% ══════════════════════════════════════════════════════════════════════════════
\section{Geometric Separation by Lift}
\label{sec:separation}
% ══════════════════════════════════════════════════════════════════════════════

\subsection{The Problem: Horizontal BLs Persist}

Horizontal velocities within the boundary layer ($z < \delta$) persist for hours:
\begin{equation}
  t_\mathrm{diff} = \frac{\delta^2}{\nu} = \frac{(2.8~\text{mm})^2}{3 \times 10^{-7}~\text{m}^2/\text{s}}
    \approx 26{,}000~\text{s} \approx 7~\text{hours}
\end{equation}

This is the molecular diffusion timescale. The horizontal flow ($u_\theta$, $u_r$)
does not decay on any operationally relevant timescale.

\subsection{Turbulent Transport Stops}

Turbulent eddies transport horizontal momentum from $z = \delta$ up to
$z = \mathrm{max\_reach}$ during their lifetime. But eddies decay rapidly:
\begin{equation}
  t_\mathrm{decay} = \frac{2\delta}{u'} = \frac{2\delta}{0.05\,V} \approx 3\text{--}14~\text{ms}
\end{equation}

Once eddies die, no further upward transport occurs. The horizontal momentum
already at max\_reach is ``frozen'' there.

\subsection{The Solution: Lift Provides Separation}

The lift moves the ion plate to $z = z_\mathrm{lift}$, placing it above max\_reach:
\begin{align}
  z_\mathrm{lift} &= \zlift~\text{mm} \\
  \mathrm{max\_reach} &= \max(\mathrm{max\_reach}_\mathrm{blade}, \mathrm{max\_reach}_\mathrm{carrier})
    = \maxreach~\text{mm} \\
  \text{Margin} &= z_\mathrm{lift} - \mathrm{max\_reach} = \margin~\text{mm} > 0
\end{align}

\subsection{Ion Trajectory}

After the lift, ions traverse from cathode ($z \approx 10$~mm) to ion plate
($z = \zlift$~mm). This path is entirely in the quiescent region ($z > \maxreach$~mm).

The ions \textbf{never encounter} the persistent horizontal boundary layers
($u_\theta$, $u_r$) which are confined to $z < \maxreach$~mm.

% ══════════════════════════════════════════════════════════════════════════════
\section{Six Mechanisms Analysis}
\label{sec:mechanisms}
% ══════════════════════════════════════════════════════════════════════════════

We analyze all potential sources of gas disturbance from MARS motion.

\subsection{Mechanism 1: Boundary Layer (Horizontal)}

\textbf{Physics:} Phases 1 and 2 create turbulent boundary layers with horizontal
velocity ($u_\theta$ from blade, $u_r$ from carrier).

\textbf{After motion stops:}
\begin{itemize}[nosep]
  \item Within BL ($z < \delta$): Horizontal velocity persists for hours
  \item Above BL ($\delta < z < \mathrm{max\_reach}$): Eddies decay in 3--14~ms
\end{itemize}

\textbf{Effect on ions: Zero} --- The lift places the ion plate above max\_reach.
Ions never traverse the horizontal flow region.

\subsection{Mechanism 2: Bulk Swirl}

\textbf{Physics:} The rotating blade transfers angular momentum to the gas,
creating bulk swirl with tangential velocity $u_\theta$.

\textbf{Key point:} Swirl velocity is purely tangential (horizontal). It has
no vertical component and cannot transport ions vertically.

\textbf{Effect on ions: Zero} --- Horizontal flow only.

\subsection{Mechanism 3: Potential (Displacement) Flow}

\textbf{Physics:} The moving blade displaces fluid, creating a potential flow
field that extends throughout the vessel.

\textbf{After motion stops:} Potential flow vanishes at the speed of sound:
\begin{equation}
  t_\mathrm{acoustic} = R/c_s = \tacoustic~\text{ms}
\end{equation}

\textbf{Effect on ions: Zero} --- Vanishes in 9~ms; ions arrive $>$1~s later.

\subsection{Mechanism 4: Pressure Pulse}

\textbf{Physics:} Blade deceleration creates a pressure pulse:
\begin{equation}
  \Delta p \sim \frac{1}{2}\rho V_\mathrm{tip}^2 \approx \dpPulse~\text{kPa}
    \quad (\sim 3\%~\text{of ambient})
\end{equation}

\textbf{After deceleration:} Pulse propagates at $c_s$ and exits vessel in $\sim$18~ms.

\textbf{Effect on ions: Zero} --- Transient; ions arrive $>$1~s later.

\subsection{Mechanism 5: Secondary Flows}

\textbf{Tip vortices:} NACA\,0012 at $\alpha = 0°$ produces no lift $\Rightarrow$
no tip vortices.

\textbf{Centrifugal effects:} Drive radial outflow, but confined to blade BL (horizontal).

\textbf{Effect on ions: Zero} --- Horizontal, confined to BL.

\subsection{Mechanism 6: Vortex Shedding}

\textbf{Physics:} NACA\,0012 at $\alpha = 0°$ is highly streamlined; shedding
is suppressed. Any shed vortices travel horizontally, confined to $z < \delta$.

\textbf{Effect on ions: Zero} --- Suppressed; horizontal if any.


% ══════════════════════════════════════════════════════════════════════════════
\section{Summary}
\label{sec:summary}
% ══════════════════════════════════════════════════════════════════════════════

\subsection{Velocity Field Summary}

The gas velocity field has three components:
\begin{align}
  u_\theta(z) &= V_\mathrm{tip} \left(\frac{z}{\delta_\mathrm{blade}}\right)^{1/7}
    && \text{(tangential, horizontal, } z < \delta_\mathrm{blade}\text{)} \\
  u_r(z) &= V_\mathrm{slide} \left(\frac{z}{\delta_\mathrm{carrier}}\right)^{1/7}
    && \text{(radial, horizontal, } z < \delta_\mathrm{carrier}\text{)} \\
  u_z(z) &= V_\mathrm{lift} \exp(-z/\delta_\mathrm{Stokes})
    && \text{(vertical, affects timing only)}
\end{align}

Only $u_\theta$ and $u_r$ can displace ions laterally. These are confined to
$z < \maxreach$~mm. The lift places the ion plate at $z = \zlift$~mm, above this region.

\subsection{Six Mechanisms: All Zero Effect}

\begin{table}[H]
\centering
\caption{Summary of gas disturbance mechanisms.}
\begin{tabular}{lllll}
\toprule
Mechanism & Component & Direction & Timescale & Effect \\
\midrule
1. Boundary Layer & $u_\theta$, $u_r$ & Horizontal & Hours & \textbf{Zero} (lift separates) \\
2. Bulk Swirl & $u_\theta$ & Horizontal & Persists & \textbf{Zero} (no $z$ component) \\
3. Potential Flow & --- & --- & 9~ms & \textbf{Zero} (vanishes) \\
4. Pressure Pulse & --- & --- & 20~ms & \textbf{Zero} (transient) \\
5. Secondary Flows & Horizontal & Horizontal & ms & \textbf{Zero} (confined to BL) \\
6. Vortex Shedding & Horizontal & Horizontal & ms & \textbf{Zero} (suppressed) \\
\bottomrule
\end{tabular}
\end{table}

\subsection{Key Performance Metrics}

\begin{table}[H]
\centering
\caption{MARS performance summary.}
\begin{tabular}{ll}
\toprule
Metric & Value \\
\midrule
Total cycle time & $t_\mathrm{total} = \ttotal$~s \\
Max reach (worst case) & \maxreach~mm \\
Lift height & $z_\mathrm{lift} = \zlift$~mm \\
Margin & \margin~mm \\
Dead zone & \deadzone~mm (\deadzonepct\% of drift) \\
Lateral ion displacement & 0 \\
\bottomrule
\end{tabular}
\end{table}

\subsection{Conclusion}

The MARS design is validated. The three-phase motion sequence completes in
\ttotal~s with all six gas disturbance mechanisms having zero effect on ion
collection. The vertical lift provides geometric separation between the ion
plate (at $z = \zlift$~mm) and the persistent horizontal boundary layers
(confined to $z < \maxreach$~mm), with a margin of \margin~mm.

\end{document}
